\documentclass[11pt]{article}
\usepackage{amsmath, amssymb}
\usepackage[margin=1in]{geometry}

\begin{document}

\begin{center}
    \textbf{\Large Homework 2}\\
    \vspace{0.2cm}
    Tashfeen Hatim Said Omran\\
    UCF ID: 5862894\\
    Class: STA4321
\end{center}

\vspace{0.5cm}

\textbf{\large Problem 2.82}

I know that $A \subset B$, meaning $A \cap B = A$. 
To find $P(B|A)$, I use the definition of conditional probability:
\[
P(B|A) = \frac{P(A \cap B)}{P(A)} = \frac{P(A)}{P(A)} = 1
\]
To find $P(A|B)$, I use the same definition:
\[
P(A|B) = \frac{P(A \cap B)}{P(B)} = \frac{P(A)}{P(B)}
\]

\textbf{\large Problem 2.88}

I am given $P(A) = 0.6$ and $P(B) = 0.3$.
\begin{itemize}
    \item[\textbf{a)}] Yes, it is possible. The probability of the union would be $P(A \cup B) = P(A) + P(B) - P(A \cap B) = 0.6 + 0.3 - 0.1 = 0.8$, which is valid since it is $\le 1$.
    \item[\textbf{b)}] I find the smallest possible value by setting $P(A \cup B) = 1$. $P(A \cap B) = P(A) + P(B) - P(A \cup B) \ge 0.6 + 0.3 - 1 = -0.1$. Since probabilities cannot be negative, the smallest possible value is $0$.
    \item[\textbf{c)}] No, it is not possible. The probability of an intersection cannot be greater than the probability of either individual event ($0.7 > 0.3$).
    \item[\textbf{d)}] The largest possible value is constrained by the smaller of the two individual probabilities. Thus, the maximum is $\min(P(A), P(B)) = 0.3$.
\end{itemize}

\textbf{\large Problem 2.93}

I am given $P(A) = 0.2$, $P(B) = 0.3$, and $P(A \cup B) = 0.4$. (Note: Assuming standard textbook complements for the missing bars in parts b and c of the prompt).
\begin{itemize}
    \item[\textbf{a)}] $P(A \cap B) = P(A) + P(B) - P(A \cup B) = 0.2 + 0.3 - 0.4 = 0.1$
    \item[\textbf{b)}] $P(\bar{A} \cup \bar{B}) = P(\overline{A \cap B}) = 1 - P(A \cap B) = 1 - 0.1 = 0.9$
    \item[\textbf{c)}] $P(\bar{A} \cap \bar{B}) = P(\overline{A \cup B}) = 1 - P(A \cup B) = 1 - 0.4 = 0.6$
    \item[\textbf{d)}] $P(A|B) = \frac{P(A \cap B)}{P(B)} = \frac{0.1}{0.3} = \frac{1}{3} \approx 0.333$
\end{itemize}

\textbf{\large Problem 2.115}

Since Team 1 passes all new cars and Team 2 rejects all new cars, a new car is rejected only if it goes to Team 2. The probability of selecting Team 2 is $0.5$. I use a binomial distribution where $X \sim \text{Binomial}(n=4, p=0.5)$ represents the number of rejected cars.
\begin{itemize}
    \item[\textbf{a)}] $P(X = 3) = \binom{4}{3} (0.5)^3 (0.5)^1 = 4(0.0625) = 0.25$
    \item[\textbf{b)}] $P(X = 0) = \binom{4}{0} (0.5)^0 (0.5)^4 = 1(0.0625) = 0.0625$
\end{itemize}

\textbf{\large Problem 2.120}

There are a total of 6 refrigerators (2 defective, 4 good). The total possible unordered combinations of assigning the 2 defective refrigerators among the 6 positions is $\binom{6}{2} = 15$. All combinations are equally likely.
\begin{itemize}
    \item[\textbf{a)}] For the last defective refrigerator to be found on the 4th test, there must be exactly 1 defective in the first 3 tests, and the 4th test must be defective. The number of such ways is $\binom{3}{1} = 3$. 
    \[
    P(\text{last D on 4th}) = \frac{3}{15} = \frac{1}{5}
    \]
    \item[\textbf{b)}] To locate both within the first 4 tests, both defective refrigerators must be placed in the first 4 positions. The number of ways is $\binom{4}{2} = 6$.
    \[
    P(\text{both in first 4}) = \frac{6}{15} = \frac{2}{5}
    \]
    \item[\textbf{c)}] I am given that exactly 1 defective is in the first 2 tests. The number of such arrangements is $\binom{2}{1}\binom{4}{1} = 8$. I need the probability that the second defective is in the 3rd or 4th position. The number of ways to have 1 defective in the first 2 tests and the remaining defective in the 3rd or 4th position is $\binom{2}{1}\binom{2}{1} = 4$.
    \[
    P(\text{2nd D in 3rd or 4th} \mid \text{exactly 1 D in first 2}) = \frac{4}{8} = \frac{1}{2}
    \]
\end{itemize}

\textbf{\large Problem 2.125}

I define $D$ as having the disease and $+$ as testing positive. I am given $P(D) = 0.01$, $P(\bar{D}) = 0.99$, $P(+ \mid D) = 0.90$, and $P(- \mid \bar{D}) = 0.90$, giving $P(+ \mid \bar{D}) = 0.10$. I use Bayes' Theorem:
\[
P(D \mid +) = \frac{P(+ \mid D)P(D)}{P(+ \mid D)P(D) + P(+ \mid \bar{D})P(\bar{D})}
\]
\[
P(D \mid +) = \frac{(0.90)(0.01)}{(0.90)(0.01) + (0.10)(0.99)} = \frac{0.009}{0.009 + 0.099} = \frac{0.009}{0.108} = \frac{1}{12} \approx 0.0833
\]
Yes, I am surprised by the low probability ($\approx 8.33\%$). I would not call this diagnostic test reliable because it produces a high proportion of false positives relative to the low disease prevalence.

\textbf{\large Problem 2.130}

I define $L$ as contracting lung cancer and $S$ as working in a shipyard. I am given $P(S \mid L) = 0.22$, $P(S \mid \bar{L}) = 0.14$, $P(L) = 0.0004$, and $P(\bar{L}) = 0.9996$. I use Bayes' Theorem:
\[
P(L \mid S) = \frac{P(S \mid L)P(L)}{P(S \mid L)P(L) + P(S \mid \bar{L})P(\bar{L})}
\]
\[
P(L \mid S) = \frac{(0.22)(0.0004)}{(0.22)(0.0004) + (0.14)(0.9996)} = \frac{0.000088}{0.000088 + 0.139944} = \frac{0.000088}{0.140032} \approx 0.000628
\]
The percentage is approximately $0.0628\%$.

\textbf{\large Problem 2.140}

I use a hypergeometric distribution where $Y$ is the number of defective refrigerators found in 3 tests. Population $N=6$, $r=2$ defective, and $n=3$ draws. The total combinations are $\binom{6}{3} = 20$.
\[
P(Y = 0) = \frac{\binom{2}{0}\binom{4}{3}}{\binom{6}{3}} = \frac{1 \times 4}{20} = \frac{1}{5}
\]
\[
P(Y = 1) = \frac{\binom{2}{1}\binom{4}{2}}{\binom{6}{3}} = \frac{2 \times 6}{20} = \frac{12}{20} = \frac{3}{5}
\]
\[
P(Y = 2) = \frac{\binom{2}{2}\binom{4}{1}}{\binom{6}{3}} = \frac{1 \times 4}{20} = \frac{1}{5}
\]

\end{document}